% !TeX root = test-main.tex

\begin{center}
{\LARGE\textbf{第1回 確認テスト}}
\end{center}

\vspace{0.4em}
\begin{flushright}
氏名：\underline{\hspace{15em}}
\end{flushright}

\vspace{0.5em}
\noindent 以下の問題は、配布テキストに沿って解答すること。

\vspace{0.5em}
\noindent\textbf{1. リスニング対策編}(5 $\times$ 10点)

\noindent\textbf{【問1】}北海道公立高校入試のリスニングにおいて、リスニングの配点は何点か。

\testblank[35点]{10em}
\vfill

\noindent\textbf{【問2】}リスニング中に聞き取れない問題があった場合、とるべき行動は何か。

\testblank[諦めて次の問題に進む]{0.92\linewidth}
\vfill

\noindent\textbf{【問3】}リスニングの問2から問4のタイミングで行うべき「次の一手」は何か。

\testblank[次の問題の設問を読んでおく]{0.92\linewidth}
\vfill

\noindent\textbf{【問4】}質問の答えを導くための重要なキーワードとは何か。

\testblank[疑問詞]{0.92\linewidth}
\vfill

\noindent\textbf{【問5】}次の慣用表現を英語に直しなさい。

「(道案内で)札幌駅までの行き方を教えていただけますか?」

\testblank[Could you tell me how to get to Sapporo Station?]{0.92\linewidth}
\vfill

\noindent\textbf{2. 英作文対策編}(5 $\times$ 10点)

\noindent\textbf{【問6】}英作文で高得点を狙うために、文法や語彙を増やすこと以上に大切な「日本語の工夫」とはどのようなことか。

\testblank[英語にしやすい日本語で表現すること（言い換え）]{0.92\linewidth}
\vfill

\noindent\textbf{【問7】}英作文の採点基準において、内容が本題に沿っていれば特に重要でないことは何か。

\testblank[内容の正しさ]{0.92\linewidth}
\vfill

\noindent\textbf{【問8】}次の日本語を、英作文しやすい表現（言い換え）に直しなさい。

「コミュニケーション能力を高めたい。」

\testblank[英語を上手に話したい。]{0.92\linewidth}
\vfill

\noindent\textbf{【問9】}問8の表現を参考に、次の日本語を、主語と動詞を含む英文1文で表現しなさい。

「コミュニケーション能力を高めたい。」

\testblank[I want to speak English well.]{0.92\linewidth}
\vfill

\noindent\textbf{【問10】}次のようにたずねられたとき、あなたはどのように答えるか。
主語と動詞を含む英文1文で自由に書きなさい。

「What do you think about introducing "World Cup" ? And why?」

\testblank[I think it's good because it is interested in soccer.]{0.92\linewidth}

\vfill
\begin{flushright}
\textbf{得点} \underline{\hspace{8em}} ／ 100点
\end{flushright}
